\section{はじめに}

ソフトウェアに含まれるOSSおよびThird-Party Library（TPL）をバイナリから特定する技術は，既知脆弱性の影響調査，ライセンス管理およびSoftware Bill of Materials（SBOM）の作成に用いられる．バイナリレベルのTPL検出では，対象バイナリの関数と既知ライブラリの関数との類似性を求め，得られた候補をライブラリやバージョンの単位に集約する．

既存手法は，命令列，制御構造，意味的特徴，関数間関係などを用いて，大量の既知関数から候補を探索する．しかし，類似性が高いことは，候補となった2関数が同じ入力に対して同じ出力を返すことを保証しない．動作の異なる候補が集約に加われば，ライブラリのバージョンを誤って特定する可能性がある．

そこで本研究は，TPL検出器が出力した関数候補を集約前に再検証する．GhidraのHigh P-codeから，既定の観測契約が可視とする戻り値，メモリ，入出力および終了に関係する式と効果を構成し，同じ論理入力に対する差の有無をSMTソルバで調べる．安全な出力差を確認できた候補はNOT\_EQとして指定バージョンの候補から除外する．差が見つからない場合でも，観測対象の経路，関数呼出しまたは副作用を十分に表現できなければEQとせずUNKNOWNとする．

評価では，EqBenchのCコードをARM64とMIPS64，ClangとGCC，O0--O3の16条件でビルドした．等価な4関数の全480ペアについて，関数全体の特徴，最適化，コンパイラおよびISAを分けて集計し，EQとUNKNOWNが分かれる要因を調べた．実High P-codeと未完了証拠を用い，ループ，戻り値格納域および関数呼出しが判定に与える影響を分析する．
